\section{Introduction}\label{sec:introduction}

Fix a monomial order $m$ on a multivariate polynomial ring.  Throughout
this article, all leading data, S-polynomials, and polynomial reductions are
taken with respect to $m$, and this dependence is suppressed from the
mathematical notation.  Let $I$ be an ideal and let $G\subseteq I$.  The set
$G$ is a Gr\"obner basis of $I$ when
\[
  \ideal{\LT(g)\mid g\in G}
  =
  \ideal{\LT(f)\mid f\in I}.
\]
The ideal on the right is the leading-term ideal of $I$.
This algebraic equality is connected with algorithmic properties of polynomial
reduction, including reduction to normal form, reduction of S-polynomials, and
the abstract Buchberger procedure
\citep{Buchberger2006,BeckerWeispfenning1993,CoxLittleOShea2025}.
A formal development must state the reduction relation and prove each
connection with this equality of ideals explicitly.

This article develops that reduction-theoretic account in Lean 4
\citep{deMouraUllrich2021}.  The development builds on Mathlib's
infrastructure for multivariate polynomials, monomial orders, leading terms,
S-polynomials, and ideals \citep{Mathlib2020}.  It also proves bridge lemmas translating finite relational reduction
sequences into finite linear-combination certificates satisfying the same
product-degree bound that appears in Mathlib's division theorem.  Its basic reduction relation allows one step to cancel any
reducible monomial in a polynomial.
Because the reduction operation is represented as a binary relation, the
standard definitions of finite reduction, joinability, equivalence closure,
normal forms, local confluence, confluence, and Church--Rosser apply without
being tied to a particular division strategy.

The formalization uses two predicates in the namespace
\lean{MonomialOrder}.  The predicate
\leandoc{Mathlib/RingTheory/MvPolynomial/GroebnerBasisCriterion}{MonomialOrder.IsGroebner}{IsGroebner}
compares the leading
terms of $G$ with the leading-term ideal of $\ideal{G}$.  The predicate
\leandoc{Mathlib/RingTheory/MvPolynomial/BuchbergerCriterion}{MonomialOrder.IsGroebnerBasis}{IsGroebnerBasis}
states that $G\subseteq I$ and compares the leading
terms of $G$ with the leading-term ideal of the specified ideal $I$.  The two
predicates are equivalent when $I=\ideal{G}$.  Throughout the article,
statements about \lean{IsGroebner} are called \emph{self-relative}, whereas
statements about \lean{IsGroebnerBasis} are called \emph{ideal-relative}.

The main contributions are as follows.
\begin{enumerate}[label=(\roman*),leftmargin=2.2em]
  \item We extend Mathlib's relation-closure API with explicit definitions of
    confluence and Church--Rosser, and we formalize normal forms, local
    confluence, Newman's lemma, and uniqueness of normal forms.
  \item We define arbitrary-term polynomial reduction and prove that its
    one-step relation admits no infinite forward reduction chain.  The proof
    uses a well-founded colexicographic order on the finite supports of
    polynomials.
  \item Assuming that every nonzero $g\in G$ has unit leading coefficient,
    we prove the exact translation
    \[
      f\eqvred{G}g
      \quad\Longleftrightarrow\quad
      f-g\in\ideal{G}.
    \]
  \item Under the same unit-leading-coefficient hypothesis, we prove that
    \[
      \ideal{\LT(f)\mid f\in\ideal{G}}
      = \ideal{\LT(g)\mid g\in G}
    \]
    is equivalent to reduction to zero for every element of $\ideal{G}$,
    Church--Rosser, local confluence, confluence, and uniqueness of normal
    forms.
  \item Assuming additionally that $R$ has no zero divisors, we prove
    Buchberger's criterion by showing that every local peak formed by two
    one-step reductions is joinable: the same-exponent case is controlled by
    an S-polynomial, while the distinct-exponent case is resolved by a direct
    reduction of the difference to zero.
  \item We formalize an abstract Buchberger procedure as a nondeterministic
    step relation.  Every step preserves the generated ideal.  After the
    unit-leading-coefficient invariant is incorporated into the state space,
    each state step strictly increases the current leading-term ideal; hence
    the state relation has no infinite forward chain over a Noetherian
    polynomial ring.  If $R$ has no zero divisors and every nonzero polynomial
    has unit leading coefficient, a terminal state is a Gr\"obner basis of
    the ideal generated by the initial polynomial set.  One step and every
    finite chain of steps preserve finiteness; consequently, in the
    finite-variable field specialization, finite input yields a finite terminal
    Gr\"obner basis.
\end{enumerate}

The phrase \emph{abstract Buchberger procedure} refers to the
nondeterministic relation defined in Section~\ref{sec:buchberger-procedure}.
It is a predicate on pairs of states: it specifies when an enlarged polynomial
set is an admissible successor of the current set, rather than defining a
function that computes one particular successor.  In particular, the
formalization does not maintain a finite queue of unprocessed critical pairs,
choose a pair by a deterministic selection function, or compute a normal form
by an executable procedure.  Under the hypotheses stated above, the state
relation admits no infinite forward chain, and every reachable terminal state
is correct; however, the formalization does not provide a function that can be
evaluated on a concrete input to return a Gr\"obner basis.
Sections~\ref{sec:related-work} and \ref{sec:future-work} compare this
relational specification with executable verified algorithms and with tactics
that obtain candidates from external computer algebra systems.

\subsection{Source organization}

The canonical Lean source is maintained in the
\lean{buchberger-formalization} branch of the author's \lean{mathlib4} fork,
since the development extends Mathlib modules directly and is intended for
upstream integration.

A separate \lean{Buchberger} repository serves as the companion artifact and
documentation site.  It does not duplicate the Lean source.  Instead, its root
module \lean{Buchberger.lean} imports the six modules from the pinned
\lean{mathlib4} snapshot.  Its Lake configuration fixes the source dependency
and Lean toolchain, while GitHub Actions verify the build and generate
\lean{doc-gen4} API documentation with declaration-level links back to the
source.

The development consists of five new files, together with additional
declarations in Mathlib's existing relation file.
Figure~\ref{fig:dependencies} shows the import dependency used in the
formalization.

\begin{figure}[t]
\centering
\setlength{\tabcolsep}{2pt}
\begin{tabular}{c}
\fbox{\strut\lean{Relation.lean}}\\[-1pt]
$\downarrow$\\[-1pt]
\fbox{\strut\lean{NormalForm.lean}}\\[-1pt]
$\downarrow$\\[-1pt]
\fbox{\strut\lean{PolynomialReductions.lean}}\\[-1pt]
$\downarrow$\\[-1pt]
\fbox{\strut\lean{GroebnerBasisCriterion.lean}}\\[-1pt]
$\downarrow$\\[-1pt]
\fbox{\strut\lean{BuchbergerCriterion.lean}}\\[-1pt]
$\downarrow$\\[-1pt]
\fbox{\strut\lean{BuchbergerAlgorithm.lean}}
\end{tabular}
\caption{Import dependency of the formal development.  The box labelled
\lean{Relation.lean} refers only to the new confluence and Church--Rosser
declarations added to that existing Mathlib file.  The other five files
constitute the reduction-theoretic Gr\"obner basis development described in
this article.}
\label{fig:dependencies}
\end{figure}

Table~\ref{tab:inventory} lists the declarations corresponding to the
principal theorems below.  In this table and
Table~\ref{tab:groebner-characterizations-lean}, \doclinkmark marks links to
the generated API documentation.

\begin{table}[t]
\caption{Selected mathematical results and their Lean declarations.}
\label{tab:inventory}
\centering
\scriptsize
\begin{tabularx}{\textwidth}{@{}>{\raggedright\arraybackslash}p{0.20\textwidth}>{\raggedright\arraybackslash}p{0.54\textwidth}>{\raggedright\arraybackslash}X@{}}
\toprule
Mathematical result & Lean declaration \doclinkmark & Source module \\
\midrule
Newman's lemma & \leandoc{Mathlib/Logic/Relation/NormalForm}{Relation.confluent_of_wellFounded_flip_of_locallyConfluent}{Relation.confluent_of_wellFounded_flip_of_locallyConfluent} & \lean{NormalForm} \\
Reduction termination & \leandoc{Mathlib/RingTheory/MvPolynomial/PolynomialReductions}{MonomialOrder.ReducesToSet.wellFounded}{MonomialOrder.ReducesToSet.wellFounded} & \lean{PolynomialReductions} \\
Finite reduction certificate with degree bound & \leandoc{Mathlib/RingTheory/MvPolynomial/PolynomialReductions}{MonomialOrder.exists_linearCombination_of_reflTransGen_with_degree_bound}{MonomialOrder.exists_linearCombination_of_reflTransGen_with_degree_bound} & \lean{PolynomialReductions} \\
Reduction equivalence and ideal congruence & \leandoc{Mathlib/RingTheory/MvPolynomial/PolynomialReductions}{MonomialOrder.eqvGen_iff_sub_mem_span}{MonomialOrder.eqvGen_iff_sub_mem_span} & \lean{PolynomialReductions} \\
Gröbner/confluence equivalence & \leandoc{Mathlib/RingTheory/MvPolynomial/GroebnerBasisCriterion}{MonomialOrder.IsGroebner.iff_confluent}{MonomialOrder.IsGroebner.iff_confluent} & \lean{GroebnerBasisCriterion} \\
Buchberger criterion & \leandoc{Mathlib/RingTheory/MvPolynomial/BuchbergerCriterion}{MonomialOrder.IsGroebner.iff_sPolynomial_reflTransGen_zero}{MonomialOrder.IsGroebner.iff_sPolynomial_reflTransGen_zero} & \lean{BuchbergerCriterion} \\
Termination of Buchberger steps & \leandoc{Mathlib/RingTheory/MvPolynomial/BuchbergerAlgorithm}{MonomialOrder.BuchbergerState.wellFounded_flip}{MonomialOrder.BuchbergerState.wellFounded_flip} & \lean{BuchbergerAlgorithm} \\
Finite terminal Gr\"obner basis & \leandoc{Mathlib/RingTheory/MvPolynomial/BuchbergerAlgorithm}{MonomialOrder.BuchbergerState.exists_finite_isGroebnerBasis_terminal_field_of_finite_variables}{MonomialOrder.BuchbergerState.exists_finite_isGroebnerBasis_terminal_field_of_finite_variables} & \lean{BuchbergerAlgorithm} \\
\bottomrule
\end{tabularx}
\end{table}
